Krátké praktické aktivity nízkého rizika, které učiteli umožní rychle vyzkoušet interakci studentů s generativními nástroji a zároveň otestovat provozní postupy (DPIA, pravidla pro sdílení dat, human‑in‑the‑loop). Před každou aktivitou si ověřte základní bezpečnostní kroky: provedení minimálního DPIA / datového plánu, vyloučení citlivých identifikátorů z odesílaných dat a jasné určení, kdo schvaluje výstupy AI (human‑in‑the‑loop) \cite{kb_malinka_chatgpt_ve_vyuce_dva_roky_pote_2024_pdf_90bc3aaf_page_12_1}. 

Přehled navrhovaných krátkých aktivit
\begin{enumerate}
  \item Icebreaker: „Co je AI?“ — Minecraft Education (30–45 min; general background suggestion)
    \begin{enumerate}
      \item Cíl: ukázat studentům konkrétní, interaktivní prostředí, kde lze demonstrovat principy učení agentů a jednoduché automatizace.
      \item Postup: krátká 10–15min demonstrační ukázka učitelem (využijte jedno z připravených světů zaměřených na AI / Hodinu kódu), poté rozdělte třídu do malých skupin, každá skupina si vyzkouší připravenou „Hour of Code“ aktivitu v Minecraft Education a společně diskutují, co v ní naznačuje učení agentů a jaké máme obavy/omezení. Minecraft Education obsahuje stovky připravených světů a některé aktivity pro Hodinu kódu jsou volně dostupné; zbytek obsahu je součástí školních licencí Microsoft 365 A3/A5 \cite{kb_ai_101_tipu_pdf_04e4a115_page_15_2}.
      \item Reflexe: rychlý záznam jedné otázky/odpovědi od každé skupiny do sdíleného dokumentu; krátká plenární diskuse o tom, co AI „umí“ a co ne.
    \end{enumerate}

  \item Praktické cvičení: „Pokrok v matematice“ — demonstration + hands‑on (45–60 min; general background suggestion)
    \begin{enumerate}
      \item Cíl: ukázat, jak nástroj může generovat cvičné příklady, sbírat postupy řešení a poskytovat předběžnou analýzu silných/slabých stránek studenta.
      \item Postup (učitel připraví zadání):
        \begin{enumerate}
          \item Učitel vybere tematickou oblast (např. zlomky, lineární rovnice) a vygeneruje sadu příkladů v nástroji typu "Pokrok v matematice".
          \item Studenti vypracují příklady a nahrají i postup řešení — nástroj předopraví odpovědi a označí oblasti, kde mohl student udělat chybu; učitel provede finální kontrolu a doplní pedagogický komentář \cite{kb_ai_101_tipu_pdf_04e4a115_page_8_1}.
          \item Diskuse: porovnejte, kde se nástroj „mýlí“ a jaká je přidaná pedagogická hodnota (diagnostika chyb, přizpůsobení obtížnosti).
        \end{enumerate}
      \item Doporučení k bezpečnosti: nepožadujte zasílání identifikovatelných osobních údajů nebo fotografií studentů do externích služeb; pokud takové data vzniknou, pseudonymizujte je před odesláním \cite{kb_malinka_chatgpt_ve_vyuce_dva_roky_pote_2024_pdf_90bc3aaf_page_12_1}.
    \end{enumerate}

  \item Dvoufázový workshop: „bez AI → s AI“ (60–90 min)
    \begin{enumerate}
      \item Cíl: pomoci studentům porozumět rozdílům v produktech lidské práce a AI‑generovaných rozšířeních, a procvičit ověřování faktů a zdrojů.
      \item Postup:
        \begin{enumerate}
          \item Fáze A (bez AI): studenti samostatně napíší krátký text (např. 300–400 slov na odborné téma nebo reflexi) během omezeného času.
          \item Fáze B (s AI): poté stejná skupina použije model (ve schválené a bezpečné instanci) k rozšíření nebo přepracování textu. Studenti dokumentují kroky, které AI provedla (prompt, varianty odpovědí, úpravy).
          \item Porovnání a reflexe: skupiny porovnají oba artefakty, hledají chyby, „halucinace“ a diskutují, jak model doplňoval informace mimo své cut‑off (ověřování aktuálnosti). Aktivita je doporučena jako způsob, jak „bezpečně startovat“ piloty a vést reflexi o ověřování a zdrojích \cite{kb_malinka_chatgpt_ve_vyuce_dva_roky_pote_2024_pdf_90bc3aaf_page_12_1,kb_ai_101_tipu_pdf_04e4a115_page_15_2}.
        \end{enumerate}
      \item Výstup: krátká skupinová prezentace o tom, kde AI přidala hodnotu a kde vznikla potřeba lidské korekce.
    \end{enumerate}

  \item Rychlá kreativní aktivita: generování výukové vizualizace (20–30 min; general background suggestion)
    \begin{enumerate}
      \item Cíl: ukázat, jak lze rychle vytvořit vizuální materiál pro hodinu (obrázek, plakát, monogram).
      \item Postup: učitel nebo skupiny studentů zformulují krátký prompt a nechají nástroj typu Designer vygenerovat návrh; následná společná úprava a sledování variant. Designer umí generovat novou grafiku, upravovat existující obrázky a umožňuje volbu poměru stran (např. širokoúhlé obrázky); mezi praktickými výstupy jsou např. monogramy, omalovánky či pozvánky \cite{kb_ai_101_tipu_pdf_04e4a115_page_51_1}.
      \item Diskuse: etické aspekty užití generované grafiky (licence, autorská práva) — krátké zamyšlení vedené učitelem.
    \end{enumerate}

  \item Krátká reflexe a deklarace použití AI (10–15 min)
    \begin{enumerate}
      \item Po každé aktivitě požádejte studenty o krátkou sebereflexi (co použili, proč, jak model měnil jejich postup) a jednoduchou deklaraci, že případné AI výstupy byly ověřeny. Transparentní oznámení a dokumentace použití AI jsou doporučenou praxí při pilotních nasazeních \cite{kb_malinka_chatgpt_ve_vyuce_dva_roky_pote_2024_pdf_90bc3aaf_page_12_1}.
      \item Vzorová formulace pro studentskou deklaraci (možno vložit do LMS):
\begin{verbatim}
Použil(a) jsem při vypracování následujících částí AI: [uveďte nástroj a stručně, jaký vstup/úpravu poskytl]. Ověřil(a) jsem fakta a upravil(a) jsem jazyk výstupu. Podpis: ...
\end{verbatim}
    \end{enumerate}
\end{enumerate}

Krátký checklist před aktivitou
\begin{itemize}
  \item Máte záznam minimálního DPIA / datového plánu pro aktivitu? (ano / ne) \cite{kb_malinka_chatgpt_ve_vyuce_dva_roky_pote_2024_pdf_90bc3aaf_page_12_1}
  \item Jsou vyloučena osobní identifikovatelná data z odesílaného obsahu? (ano / ne) \cite{kb_malinka_chatgpt_ve_vyuce_dva_roky_pote_2024_pdf_90bc3aaf_page_12_1}
  \item Je jasně určeno, kdo finálně kontroluje a schvaluje výstupy AI (human‑in‑the‑loop)? (ano / ne) \cite{kb_malinka_chatgpt_ve_vyuce_dva_roky_pote_2024_pdf_90bc3aaf_page_12_1}
  \item Má učitel připravenou krátkou reflexní otázku pro studenty (záznam o ověření výstupů)? (ano / ne)
\end{itemize}

Poznámky a doporučení
\begin{itemize}
  \item Začněte s nehodnoticím obsahem (demonstrační aktivity, cvičné listy, vizuály), kde případné chyby lze snadno opravit — to snižuje riziko negativního dopadu na studenty (general background knowledge).
  \item Po ukončení pilotu zaznamenejte artefakty (studentské výstupy, prompty, záznamy o opravách) a krátce vyhodnoťte pedagogický přínos a provozní rizika; výsledky slouží pro další úpravy a případné širší nasazení.
\end{itemize}